\documentclass[11pt,letterpaper]{article}
\usepackage[margin=1in]{geometry}\usepackage{amsmath,amssymb,booktabs,array,microtype}\usepackage[T1]{fontenc}\usepackage{lmodern}
\usepackage[colorlinks=true,linkcolor=blue,citecolor=blue,urlcolor=blue]{hyperref}
\setlength{\parindent}{0pt}\setlength{\parskip}{5pt}\sloppy
\title{What a Clock Can and Cannot Do:\\ a Cold Dark Sector Derived, a Force Law Excluded,\\ and the Rate of Cosmic Time Certified}
\author{Carl P.\ Zimmerman\\ \small Briar Creek Tech\\ \small AI-assisted research programme; not peer reviewed}
\date{12 September 2026}
\begin{document}\maketitle
\begin{abstract}
A cuscuton-clock sector with a projected-gradient scalar has been shown to produce a cold, clustering component by driving itself onto a critical surface, and to pass nine successive cosmological and solar-system gates \cite{paper19,paper20,paper21,paper22}. Every one of those derivations set the cubic coupling $\gamma$ to zero. We restore it and report two things, one negative and decisive, one positive and certified. Negative: in the static spherically symmetric limit the action's two gradient operators give first integrals $r^2(\chi')^3$ and $\gamma r(\chi')^2$, hence rotation curves rising as $r^{1/6}$ and $r^{1/4}$; flat curves require a first integral $r^2(\chi')^2$, that is $F\propto Y^{3/2}$, the AQUAL deep-MOND kinetic term. Both available exponents are strictly positive, so no value of $\gamma$ and no mixture of the two operators reaches zero. \emph{This action cannot produce flat rotation curves.} A partial consolation is that its two terms cancel exactly at zero gradient, so the deep-MOND boundary condition $\mu\to0$ holds without being imposed; only the exponent is wrong. Positive: restoring $\gamma$ leaves the energy density and the clock's charge untouched --- the cubic cancels out of both --- and survives only in the pressure, correcting the clock identity to $s_0-1=w/m_{\rm rel}-2\gamma q^2\dot q/U$, a correction three parts in a million of the leading term at the coupling used. The identity itself, and its consequence that the clock exceeds proper time if and only if the sector's equation of state is positive, are certified in Lean~4. The honest scope is therefore: this construction derives a cold dark sector and the rate of cosmic time from a clock, and does not derive the force law. That is a narrower claim than a theory of gravity and a sharper specification for one.
\end{abstract}

\section{The question that was assumed past}
The action \cite{paper13} is, with $c=1$ and signature $(-{+}{+}{+})$,
\begin{equation}
S=\int d^4x\sqrt{-g}\Big[\tfrac{M^2}{2}(R-2\Lambda)+P(X,\tau)-V(\tau)+sW(Y,\tau)+\gamma X\Box\chi\Big]+S_r+S_b,
\label{action}
\end{equation}
with $s=\sqrt{-\nabla_\mu\tau\nabla^\mu\tau}$, $n_\mu=-\partial_\mu\tau/s$, $Q=n^\mu\partial_\mu\chi$, $Y=(g^{\mu\nu}+n^\mu n^\nu)\partial_\mu\chi\partial_\nu\chi$, $X=Q^2-Y$. Its cosmological reduction and the gate results built on it \cite{paper19,paper20,paper21,paper22} all took $\gamma\to0$, which is harmless for a homogeneous background. It is not harmless for whether (\ref{action}) produces a modified force law, because $\gamma X\Box\chi$ is the only Galileon-type operator present.

\section{The static limit with $\gamma=0$: right boundary condition, wrong exponent}
With the clock at rest, $Q=0$ and $X=-Y$, so the entire $Y$-dependence is $F(Y)=P(-Y)+W(Y)$, and for such a Lagrangian the field equation is $\nabla\!\cdot\!\big(2F'(Y)\nabla\chi\big)=\text{source}$: the quantity $2F'$ plays exactly the role of the interpolating function. With the closure's $P$ and $W$,
\begin{equation}
2F'(Y)=\frac{2d}{\sqrt{1+Y/\ell}}-\frac{2Ud}{U+2dY},
\end{equation}
whose two terms are $+2d$ and $-2d$ at $Y=0$. They cancel \emph{exactly}:
\begin{equation}
2F'(0)=0 .
\end{equation}
That is the deep-MOND boundary condition $\mu\to0$ as the acceleration vanishes, and it holds without being imposed. Expanding,
\begin{equation}
2F'(Y)=\Big(\frac{4d^2}{U}-\frac{d}{\ell}\Big)Y+O(Y^2),
\end{equation}
linear in $Y$ and therefore \emph{quadratic} in the gradient, where MOND requires the first power. A by-product is a health condition the cosmological work could not see, because it worked at $Y=0$: the coefficient is positive only for $4d\ell>U$.

\section{Restoring the cubic}
\emph{Background.} Carrying $\gamma$ through the sector's stress \cite{paper13}, the explicit cubic piece of $P$ contributes $+6\gamma Hq^3$ to $2q^2P_X$ while the stress carries $-6\gamma Hq^3$; the same cancellation occurs in the current. Hence
\begin{equation}
\rho=\frac{U}{m_{\rm rel}},\qquad j=\frac{2qd}{m_{\rm rel}}\qquad\text{unchanged},\qquad p=U(s_0-1)+2\gamma q^2\dot q ,
\end{equation}
so the cubic survives in the pressure alone and the identity of \cite{paper21} becomes
\begin{equation}
\boxed{\ s_0-1=\frac{w}{m_{\rm rel}}-\frac{2\gamma q^2\dot q}{U}\ }
\label{identity}
\end{equation}
At the coupling used in the candidate's own runs, $\gamma=10^{-6}$, the correction is $6\times10^{-10}$ against a leading term of $2\times10^{-4}$, a ratio of $3\times10^{-6}$: the closed-form family and the acoustic-scale bound of \cite{paper21} stand exactly as derived.

\emph{Static limit.} In spherical symmetry the cubic is a cubic Galileon, contributing a first-integral term $\propto\gamma r(\chi')^2$ alongside $r^2\,2F'(Y)\chi'$. Setting each equal to a constant and reading off $v^2=r|g|$:
\begin{table}[h]\centering\caption{What each operator gives.}\label{tab}
\begin{tabular}{lcc}\toprule
operator & first integral & rotation curve\\\midrule
$P$ and $W$ together, $F'\propto Y$ & $r^2(\chi')^3$ & $v\propto r^{1/6}$, rising\\
cubic Galileon & $\gamma\,r(\chi')^2$ & $v\propto r^{1/4}$, rising faster\\
what flat curves require, $F\propto Y^{3/2}$ & $r^2(\chi')^2$ & $v\propto r^{0}$, flat\\\bottomrule
\end{tabular}\end{table}

\section{The no-go}
Both available exponents are strictly positive, $1/6$ and $1/4$. A sum of first-integral terms that each force a rising curve yields an exponent between them, and $0$ is not in $[1/6,1/4]$. Therefore \emph{(\ref{action}) cannot produce flat rotation curves for any value of $\gamma$}. What flat curves require is a first integral linear in the gradient, that is $F\propto Y^{3/2}$; the action's two $Y$-dependent pieces are a logarithm and a square root, and neither reduces to a three-halves power in any limit. The operator is not present.

This is a statement about the static limit of one action, not about the framework's kernel, which is an independent ingredient and is unaffected.

\section{How cosmic time runs, certified}
The positive content is (\ref{identity}) and what follows from it, now machine-checked in Lean~4 (\texttt{Mondlean.lean}, 114 theorems, no \texttt{sorry}, axioms \texttt{propext}, \texttt{Classical.choice}, \texttt{Quot.sound}):
\texttt{clock\_rate\_from\_conservation} --- imposing only $w=p/\rho$ on the forms above yields (\ref{identity}): the clock's rate relative to proper time is not an input, it is fixed by the sector's own pressure.
\texttt{clock\_rate\_is\_pressure\_over\_margin} --- with $\gamma=0$, $s_0-1=w/m_{\rm rel}$ exactly.
\texttt{clock\_runs\_fast\_iff\_pressure\_positive} --- with a positive margin, $s_0>1$ if and only if $w>0$.
Chained with the stability equivalence of \cite{paper19}, which makes the sector gradient-unstable precisely when $s_0>1$, and with the criticality that instability drives: \emph{a positive dark-sector pressure makes cosmic time run fast, which destabilises the sector, which its own nonlinearity cures at a critical gradient whose marginal state is exactly cold matter.} Every link is certified algebra; the only free number is $w$, bounded above at $10^{-4}$ by the acoustic scale \cite{paper21} and required strictly positive by the criticality.

\section{The honest scope, and what it specifies}
This construction derives a cold, clustering dark sector and the rate of cosmic time from a clock. It does not derive the force law. The nine gates passed in \cite{paper19,paper20,paper21,paper22} are dark-matter gates, and the sector passes them as a dark component.

That is also a specification. Any action intended to do both jobs must carry a $Y^{3/2}$ operator alongside the clock, and the derived family fixes the rest: $U\propto a^{-3(1+w)}$, $d\propto a^{-3(1-w)}$, $q\propto a^{-3w}$, with $s_0-1=w/m_{\rm rel}$, $4d\ell>U$, and $0<w\lesssim10^{-4}$. That is a far narrower target than the one this programme began with.

\section{What is not claimed}
Static spherical weak field with the clock at rest; the matter coupling to $\chi$ is taken to be the standard scalar-tensor form, so identifying the first integral with the enclosed mass assumes it; the Galileon's first integral is read off its variation rather than solved in full; the large-gradient branch of $F$ is untreated and carries no standard static solution, since there $F'$ falls as $1/x$. Lean certifies the algebra and the logic of section~6; that $\rho$ and $p$ take the stated forms follows from varying (\ref{action}) and is not formalised. And $\kappa=\tfrac12$ remains fitted and provably underivable by this class of actions \cite{repo}.

\section{Reproducibility}
\texttt{fable\_independent\_2026/L205\_does\_the\_family\_give\_mond.py} and \texttt{L206\_cubic\_restored.py} with their committed outputs (commits \texttt{ac2052ae0}, \texttt{da0fafdf5}); the Lean file \texttt{fable\_independent\_2026/lean\_2026/Mondlean.lean} (commit \texttt{0eee1a513}).

\begin{thebibliography}{9}
\bibitem{repo} The research repository, \url{https://github.com/carlzimmerman/zimmerman-formula} (2026).
\bibitem{paper13} C.~P.~Zimmerman, \emph{A Primordial-Clock Relativistic MOND Candidate: Status Report}, doi:10.5281/zenodo.22699828 (2026).
\bibitem{paper19} C.~P.~Zimmerman, \emph{Gradient-Driven Criticality}, doi:10.5281/zenodo.22727111 (2026).
\bibitem{paper20} C.~P.~Zimmerman, \emph{The Coldness Is Dynamical}, doi:10.5281/zenodo.22727860 (2026).
\bibitem{paper21} C.~P.~Zimmerman, \emph{The Clock Rate Is the Dark Matter Pressure}, doi:10.5281/zenodo.22728545 (2026).
\bibitem{paper22} C.~P.~Zimmerman, \emph{A Clock That Cannot Be Dragged}, doi:10.5281/zenodo.22728789 (2026).
\bibitem{bekenstein} J.~Bekenstein, M.~Milgrom, Astrophys.\ J.\ \textbf{286}, 7 (1984).
\bibitem{cuscuton} N.~Afshordi, D.~J.~H.~Chung, G.~Geshnizjani, Phys.\ Rev.\ D \textbf{75}, 083513 (2007).
\end{thebibliography}
\end{document}
